\documentclass{article}
\usepackage{tcee-repaso}
\tcetema[Relaciones exteriores y cooperación al desarrollo de la UE]{3.B.42}{Las relaciones económicas exteriores de la Unión Europea. La política de cooperación al desarrollo de la Unión Europea.}
\tcefecha{19/11/2026}
\begin{document}
\cabecerarepaso
\begin{columnas}

% ==========================================================================
\bloque{Esquema del cante}
\esquemacante{%
\apartadocante{Introducción: «más vale vecino cercano que hermano lejano»}{2,5}
\apartadocante{I. Visión global: peso, marco y organizaciones}{8,5}
\subcante{1. Peso económico (el euro)}{3}
\subcante{2. Justificación y marco jurídico e institucional}{3}
\subcante{3. Organizaciones y foros}{2,5}
\apartadocante{II. Visión regional: vecinos, socios y rivales}{10}
\subcante{1. Ampliación, vecindad, sanciones (tope)}{6}
\subcante{2. Factores, grandes socios y retos}{4}
\apartadocante{III. La cooperación al desarrollo}{7,5}
\subcante{1. Consideraciones generales}{2}
\subcante{2. La UE y la AOD (fórmula común)}{5,5}
\apartadocante{Conclusión}{1,5}}
\emph{Gráficos del cante}: el euro en las reservas (\hyperref[rep:I]{I}), círculos de las relaciones exteriores (\hyperref[rep:II]{II}), AOD por Estados e IVDCI (\hyperref[rep:III]{III}). Modelos en el anexo: efecto Bruselas $F>c\,M_R$, condicionalidad $C\leq g+p\,\delta^{T}B$ y reparto de la ayuda $m_1=m_2$.

% ==========================================================================
\bloque[2,5]{Introducción}
\begin{itemize}
\item \textbf{Enganche}: «Más vale vecino cercano que hermano lejano». En 2025 EE.\,UU. puso un 15\,\% de arancel a la UE y recortó su ayuda a menos de la mitad; la UE siguió siendo el primer donante, lanzó el Pacto por el Mediterráneo y abrió en 2026 las negociaciones con Ucrania y Moldavia. ¿Para qué usa la UE su peso económico fuera, y le basta con comerciar?
\item \textbf{Relevancia}: 2.ª economía del mundo (17,9\,\% del PIB), 2.º exportador e importador de bienes, 1.º de servicios, 1.er socio de unos 80 países; el \concepto{Equipo Europa} da la mitad de la AOD del CAD. España: Mediterráneo, América Latina \vertema{4.A.26}.
\item \textbf{Contextualización}: \concepto{relaciones económicas exteriores}: un conjunto de políticas (comercial, ampliación, vecindad, desarrollo, humanitaria, macrofinanciera, sanciones). Comercio exclusivo; desarrollo y humanitaria \emph{paralelas}; PESC por unanimidad \vertema{3.B.36}. Comercio \vertema{3.B.41}; financiación del desarrollo \vertema{3.B.29}.
\item \textbf{Problemática}: ¿cuánto pesa, por qué en común y con qué reglas y actores? ¿Por qué trata distinto a candidatos, vecinos y socios? ¿Qué hace por el desarrollo y con qué resultados?
\item \textbf{Estructura}: \hyperref[rep:I]{I} visión global; \hyperref[rep:II]{II} visión regional; \hyperref[rep:III]{III} cooperación al desarrollo.
\end{itemize}

% ==========================================================================
\bloque[8,5]{I. Visión global\label{rep:I}}
\begin{ideaclave}
La UE pesa en comercio, inversión, ayuda y moneda, y actúa fuera con muchas manos y casi solo con instrumentos económicos: mercado, normas, dinero y moneda. Ese es su poder y su límite.
\end{ideaclave}

\sub{I.1. Peso económico\label{rep:I.1}}
\begin{itemize}
\item \emph{Comercio e inversión}: 5,5\,\% de la población, 17,9\,\% del PIB (14,0\,\% en PPA). Bienes: 2.º exportador (13,8\,\%) e importador (12,9\,\%); servicios: 1.º (22,7\,\%). EE.\,UU. 1.er cliente (superávit $\approx$197.000 M€), China 1.er proveedor (déficit $\approx$375.000 M€). IED: 9,3 billones fuera, 7,0 dentro. \emph{Ayuda}: la mitad de la AOD del CAD.
\item \emph{Moneda internacional}: 3 funciones fuera (privados: facturar, pagar, ahorrar; gobiernos: ancla, intervención, reservas). \emph{Determinantes}: peso, confianza, externalidades de red, mercados profundos y activo seguro (el punto débil del euro). \emph{Ventajas}: señoreaje, menos riesgo cambiario, autonomía frente a sanciones. \emph{Inconvenientes}: volatilidad, apreciación en crisis, carga para el BCE.
\item \emph{Indicadores 2025 (BCE, junio de 2026)}: 2.ª moneda, índice $\approx$20\,\%; reservas 20,2\,\% (dólar 57\,\%; oro 27\,\% de las reservas oficiales); deuda 24,6\,\%; SWIFT $\approx$23\,\%; divisas 28,5\,\%; facturación 60,0\,\% de las exportaciones y 53,1\,\% de las importaciones. Máximo de reservas 24\,\% (2009); $\approx$20\,\% desde 2015. Ancla de una treintena de países (franco CFA, Balcanes); Bulgaria, 21.º miembro (2026). «Momento del euro global» (\autor{Lagarde}, 2025); EUREP global (2026) \vertema{3.B.43}.
\end{itemize}

\pareja{euro-reservas}{El euro, en torno al 20\,\% de las reservas desde 2015; el dólar baja}{comercio-socios}{EE.\,UU., China y Reino Unido, los 3 grandes socios (2025)}

\sub{I.2. Justificación y marco\label{rep:I.2}}
\begin{itemize}
\item \emph{Teoría}: federalismo fiscal (\autor{Oates}, 1972; \autor{Alesina}, \autor{Angeloni} y \autor{Etro}, 2005): centralizar si escala y externalidades $>$ heterogeneidad. A favor: la unión aduanera lo exige, poder de negociación, poder normativo (\concepto{efecto Bruselas}, \autor{Bradford}, 2012: anexo), bienes públicos, escala. En contra: preferencias distintas $\Rightarrow$ comercio exclusivo, desarrollo paralelo, PESC unánime.
\item \emph{Base}: arts. 21-22 TUE; TFUE, 5.ª parte (205-222): 206-207 comercio; 208-211 desarrollo; 212-213 terceros y macrofinanciera; 214 humanitaria; 215 sanciones; 216-219 acuerdos; 220-221 organizaciones y delegaciones; 222 solidaridad.
\item \emph{Actores y decisiones}: Consejo Europeo (\autor{Costa}); Consejo de Asuntos Exteriores (AR); AR/VP \autor{Kallas} y SEAE (145 delegaciones); Comisión (TRADE, INTPA, ENEST y MENA, ECHO, ECFIN); PE. Comercio y desarrollo: PLO y mayoría cualificada; PESC: unanimidad; acuerdos: art. 218.
\item \emph{Evolución}: Roma (1957; arancel común 1968); CPE (1970); Maastricht (PESC, 1992); AR (1999); Lisboa (2009: personalidad jurídica, AR/VP, SEAE); giro geoeconómico: rival sistémico (2019), Global Gateway (2021), seguridad económica (2023).
\item \emph{Instrumentos}: acuerdos (asociación, comerciales, asociación económica, cooperación; \emph{mixtos}: Dictamen 2/15); rúbrica 6 (113.700 M€: IVDCI 79.462, IPA III 14.162, humanitaria 11.569) \vertema{3.B.37}; Mecanismo para Ucrania (50.000); macrofinanciera; sanciones.
\end{itemize}

\sub{I.3. Organizaciones y foros\label{rep:I.3}}
\begin{itemize}
\item \emph{Base}: art. 220 TFUE (ONU, OCDE…); 221 delegaciones; art. 34 TUE (Consejo de Seguridad); art. 138 TFUE (posiciones del euro en el FMI).
\item \emph{Miembro}: OMC (1995; habla la Comisión; MPIA 2020; MC14 Yaundé 2026 sin moratoria electrónica), FAO (1991), G-20 (1999), BERD (UE y BEI, 3\,\% cada uno). \emph{Miembro no enumerado}: G-7 (1977). \emph{Observador}: ONU (reforzado, Res. 65/276, 2011), FMI (el BCE, 1999; EURIMF; los Estados $\approx$1/4 de los votos sin silla única) \vertema{3.B.22}. \emph{OCDE}: la Comisión sin voto; miembro del CAD; 22 Estados \vertema{3.B.19}. \emph{FSB}: Comisión y BCE; \emph{BPI}: BCE. Banco Mundial: los Estados \vertema{3.B.30}.
\end{itemize}

% ==========================================================================
\bloque[10]{II. Visión regional\label{rep:II}}
\begin{ideaclave}
Círculos concéntricos: cuanto más cerca, más integración y más condiciones. Candidatos: adhesión; vecinos: mercado y dinero por reformas; socios: acuerdos o sanciones.
\end{ideaclave}

\graficorepaso[\linewidth]{circulos-relaciones}{UE-27, candidatos, vecindad, grandes socios y países en desarrollo}

\sub{II.1. La economía como instrumento geoestratégico\label{rep:II.1}}
\begin{itemize}
\item \emph{Ampliación}: art. 49; \concepto{criterios de Copenhague} (1993: políticos, económicos, acervo) y Madrid (1995: administración). Metodología 2020: 6 bloques, fundamentos primero y últimos, reversible. IPA III 14.162 M€; Plan de Crecimiento de los Balcanes 6.000 (2.000 + 4.000 en préstamos); Moldavia 1.900; Mecanismo para Ucrania 50.000 ($>$32.000 pagados) y Préstamo de Apoyo 90.000 (2026-2027).
\item \emph{Situación (paquete del 4-11-2026)}: 9 candidatos + Kosovo; Montenegro 21/33 capítulos cerrados; Albania 6 bloques abiertos (fin: 2027); Ucrania y Moldavia, 4 bloques (fin: 2028); hojas de ruta e integración gradual; Serbia sin avances; Georgia candidata «de nombre»; Turquía parada desde 2018.
\end{itemize}

\graficorepaso[\linewidth]{etapas-adhesion}{De la solicitud (art. 49) a la ratificación por todos}

\begin{itemize}
\item \emph{Condicionalidad} (\autor{Schimmelfennig} y \autor{Sedelmeier}, 2004; anexo): reforma si la adhesión, más las recompensas intermedias, cubre el coste interno; funciona si es grande, creíble y cercana (2004 sí; Balcanes tras Salónica, no); fatiga; Plan de Crecimiento = pagar antes; tras entrar, condicionalidad del Estado de Derecho (2020).
\end{itemize}

\begin{itemize}
\item \emph{Vecindad} (2004; art. 8 TUE): «más por más». \emph{Este}: Asociación Oriental (2009; Bielorrusia suspendida, 2021); ZLCAP con Georgia, Moldavia y Ucrania (2014). \emph{Sur}: 10 países; Unión por el Mediterráneo (2008, 43 miembros, Barcelona). DG MENA (2025); \concepto{Pacto por el Mediterráneo} (16-10-2025: personas, economías, seguridad y migración); asociaciones con Túnez (2023), Mauritania (2024), Egipto (2024, 7.400 M€), Jordania (2025, 3.000 M€). Condicionalidad débil, más transaccional.
\item \emph{PESC y sanciones}: decisión unánime (art. 29 TUE) + reglamento (art. 215 TFUE). Rusia: 21 paquetes (último 23-7-2026); $\approx$210.000 M€ del banco central inmovilizados sin plazo (dic. 2025); fin del gas ruso (Reglamento 2026/261: GNL largo plazo 1-1-2027, gasoducto 30-9-2027); $>$550 buques de la flota en la sombra.
\end{itemize}

\begin{itemize}
\item \emph{Tope al petróleo} (G-7 y UE, 2022; \autor{Johnson}, \autor{Rachel} y \autor{Wolfram}, 2023): quitar renta a Rusia sin quitar su crudo del mercado; servicios occidentales solo bajo el tope; renta $R=[\theta\bar P+(1-\theta)(P^*-k)-c]\bar Q$; punto débil, la flota en la sombra. Tope dinámico: 60 \$ $\to$ 47,6 (sept. 2025) $\to$ 44,1 (feb. 2026, congelado hasta julio de 2027). Descuento del Urales de 32 \$ (marzo de 2023); ingresos rusos del petróleo y el gas: 35,8\,\% $\to$ 22,7\,\% (2021-2025). Contención y señal más que coerción.
\end{itemize}

\graficorepaso[0.8\linewidth]{tope-petroleo}{La renta entre $\bar P$ y $P^*$ pasa a los compradores}

\sub{II.2. Factores, grandes socios y retos\label{rep:II.2}}
\begin{itemize}
\item \emph{Factor a factor}: bienes (art. 207, \vertema{3.B.41}); servicios; capitales (art. 63, también con terceros; control de IED, Reglamento 2019/452); personas (arts. 77-80; Pacto de Migración y Asilo, 2024, aplicable desde junio de 2026).
\item \emph{Avanzados}: EE.\,UU. (sin acuerdo; Turnberry y declaración del 21-8-2025: arancel máximo del 15\,\%; el Supremo anula los aranceles de emergencia, 2026); Reino Unido (ACC 2021; cumbres 19-5-2025 y 20-11-2026); EEE y Suiza \vertema{3.B.35}; Japón, Corea, Canadá.
\item \emph{Rivales}: China («socio, competidor y rival sistémico», 2019; \emph{de-risking}; déficit $\approx$375.000 M€); Rusia (sanciones).
\item \emph{Emergentes}: CELAC (Santa Marta 2025, 45.000 M€), Mercosur (firmado 2026, provisional); UA-UE (Luanda 2025, 150.000 M€); India (concluido 2026); CCG (1.ª cumbre 2024).
\item \emph{Retos}: mundo de bloques; una sola voz (mayoría cualificada, art. 31.3 TUE); ampliación; marco 2028-2034; euro global.
\end{itemize}

% ==========================================================================
\bloque[7,5]{III. La cooperación al desarrollo\label{rep:III}}
\begin{ideaclave}
1.er donante del mundo, pero de las donaciones a las asociaciones e inversión (Global Gateway), con más interés propio y con recortes en varios Estados.
\end{ideaclave}

\sub{III.1. Consideraciones generales\label{rep:III.1}}
\begin{itemize}
\item \concepto{AOD}, 4 requisitos del CAD: (1) a países de la lista o multilaterales; (2) oficial; (3) objetivo principal, el desarrollo; (4) concesional: donación equivalente desde 2018, elemento de donación del 45/15/10\,\% (descuento del 9/7/6\,\%) \vertema{3.B.29}. Tipos: financiera (apoyo presupuestario), técnica, económica, humanitaria. 0,7\,\% (ONU, 1970).
\item \emph{En el mundo}: CAD 2025: 174.300 M\$, 0,26\,\% ($-$23,1\,\%, la mayor caída). ODM $\to$ ODS; Monterrey, Doha, Addis Abeba, \emph{Compromiso de Sevilla} (2025); París (2005), Accra, Busan.
\end{itemize}

\graficorepaso[0.8\linewidth]{aod-ue-estados}{Solo LU, SE y DK superan el 0,7\,\%; España, 0,27\,\%}

\sub{III.2. La UE y la AOD\label{rep:III.2}}
\begin{itemize}
\item \emph{Datos 2025}: los 27, $\approx$87.000 M\$ (0,41\,\% de la RNB, la mitad del CAD); instituciones, 23.100 M€ ($-$13,8\,\%); miembros UE del CAD 0,47 $\to$ 0,42\,\%; Alemania 1.er donante (0,68 $\to$ 0,56\,\%); recortes en BE, FI, FR, DE, NL, SE.
\item \emph{Justificación}: pobreza (capital, trampas, instituciones: \autor{Acemoglu}, \autor{Johnson} y \autor{Robinson}, Nobel 2024); por qué la UE: equidad, eficiencia, geopolítica, historia, escala y coordinación; a quién: selectividad (\autor{Collier} y \autor{Dollar}, 2002: $m_1=m_2$, de 30 a 80 millones al año; anexo). \emph{Competencias}: paralela (art. 4.4). \emph{Marco}: arts. 21 TUE y 208 TFUE (pobreza), \concepto{coherencia de políticas}, Consenso (2017: 5 P). \emph{Decisiones}: PLO (209), Comisión (INTPA, ECHO), coordinación (210); humanitaria (214).
\item \emph{Evolución}: Roma (1957) y 1.er FED (1959); Yaundé (1963, 1969); Lomé (1975-2000: no recíprocas, STABEX); SPG (1971); Maastricht (1992); Cotonú (2000: AAE recíprocos); Todo menos armas (2001); \concepto{Samoa} (15-11-2023; provisional desde 1-1-2024; base común + 3 protocolos; aún sin vigor).
\end{itemize}

\graficorepaso[\linewidth]{cronologia-cooperacion}{De la asociación con las excolonias a las asociaciones e inversión}

\begin{itemize}
\item \emph{Articulación}: \concepto{IVDCI-Europa Global} (79.462 M€; geográficos 60.388: África subsahariana $\geq$29.181, vecindad $\geq$19.323, Asia 8.489, América 3.395; temáticos 6.358; respuesta rápida 3.182; colchón 9.534; 93\,\% AOD, 30\,\% clima, 20\,\% desarrollo humano); \concepto{FEDS+} (garantía hasta 53.449 M€; hasta 232.000 de inversión); \concepto{Equipo Europa} (2020); \concepto{Global Gateway} (2021: 300.000 M€; 306.000 movilizados en 2021-2024); ECHO (1.900 M€ iniciales en 2026); SPG y AAE \vertema{3.B.41}.
\end{itemize}

\graficorepaso[0.9\linewidth]{ivdci}{África subsahariana y vecindad, 6 de cada 10 euros del IVDCI}

\begin{itemize}
\item \emph{Retos}: \emph{Europa Global} 2028-2034: $\approx$200.000 M€ (0,14\,\% de la RNB), fusiona IVDCI, IPA y humanitaria; AOD mínima del 90\,\%; sin suelos temáticos; Ucrania aparte \vertema{3.B.37}.
\item \emph{España}: 4.554 M€ (0,27\,\%) en 2025; Ley 1/2023: 0,7\,\% en 2030, 10\,\% humanitaria, FEDES (sustituye a FONPRODE); Plan Director 2024-2027; AECID en la cooperación delegada \vertema{4.A.26} \vertema{4.A.24}.
\end{itemize}

% ==========================================================================
\bloque[1,5]{Conclusión}
\begin{itemize}
\item \textbf{Recapitulación}: En conclusión, políticas con base común y poder económico (mercado, normas, dinero y euro); círculos: adhesión condicionada, vecindad, acuerdos y sanciones; 1.er donante que pasa de la ayuda a la inversión.
\item \textbf{Relevancia}: actor o escenario en un mundo de bloques; la vía de España hacia el Mediterráneo, América Latina y África.
\item \textbf{Extensiones}: \vertema{3.B.41} \vertema{3.B.29} \vertema{3.B.30} \vertema{3.B.37} \vertema{3.B.43} \vertema{4.A.26} \vertema{4.A.24}.
\item \textbf{Opinión}: más poder del que usa y menos del que cree; la ampliación es su mejor política exterior, si la promesa es creíble y la UE se reforma.
\item \textbf{Idea final}: más vale vecino cercano que hermano lejano: cuando el hermano lejano sube aranceles y recorta ayuda, la fuerza de la UE está en vecinos y socios; su peso económico le sirve para más que comerciar, si cumple lo que promete.
\end{itemize}

% ==========================================================================
\bloque{Anexo: 3 modelos}
\begin{itemize}
\item \textbf{Efecto Bruselas} (\autor{Bradford}, 2012). \emph{Idea}: exportación unilateral de normas por el mercado (poder normativo; poder blando de \autor{Nye}). \emph{Supuestos}: ventas $M_{UE}$ y $M_R$; coste de la norma $c$ por unidad; 2 líneas cuestan $F$; objetivo poco móvil. \emph{Desarrollo}: una sola línea si
\begin{formula}
\funcionclave{\text{efecto Bruselas}\iff F>c\,M_R}
\end{formula}
\emph{Implicaciones}: 5 condiciones (mercado grande, capacidad reguladora, normas estrictas, objetivos poco móviles, indivisibilidad); de hecho $\to$ de derecho (RGPD). \emph{Evidencia}: RGPD, REACH. \emph{Valoración}: depende del tamaño del mercado; riesgo para la innovación (\emph{Draghi}).
\item \textbf{Condicionalidad} (\autor{Schimmelfennig} y \autor{Sedelmeier}, 2004). \emph{Supuestos}: coste $C$; adhesión de valor $B$ con probabilidad $p$ en $T$ años; descuento $\delta$; recompensas intermedias $g$. \emph{Desarrollo}: reforma si
\begin{formula}
\funcionclave{C\leq g+p\,\delta^{T}B}
\end{formula}
\emph{Extensiones}: inconsistencia temporal tras entrar; geopolítica. \emph{Evidencia}: \autor{Vachudova} (2005); \autor{Campos}, \autor{Coricelli} y \autor{Moretti} (2019): $\approx$+10\,\% de renta en 10 años. \emph{Valoración}: actor único; $p$ depende de la unanimidad.
\item \textbf{Reparto eficiente de la ayuda} (\autor{Collier} y \autor{Dollar}, 2002). \emph{Supuestos}: $A_1+A_2=\bar A$; $m_i$ decreciente, mayor con más pobres y mejores políticas. \emph{Desarrollo}:
\begin{formula}
\funcionclave{m_1(A_1^*)=m_2(A_2^*)}
\end{formula}
\emph{Extensiones}: coordinación entre donantes (\autor{Knack} y \autor{Rahman}; \autor{Bigsten} y \autor{Tengstam}). \emph{Evidencia}: no robusto (\autor{Easterly}, \autor{Levine} y \autor{Roodman}, 2004); reparto por lazos (\autor{Alesina} y \autor{Dollar}, 2000). \emph{Valoración}: castiga a los frágiles; \autor{Easterly}: buscadores.
\end{itemize}

\pareja{efecto-bruselas}{Por encima de la diagonal, una sola línea con la norma UE}{condicionalidad}{$X$ reforma con $T_1$, no con $T_2$, y otra vez con $g$}

\graficorepaso[0.8\linewidth]{asignacion-ayuda}{$E$: $m_1=m_2$; el reparto igual pierde el triángulo}

% ==========================================================================
\bloque{Autores y fechas}
\begin{itemize}
\item 1957 Roma; 1972 \autor{Oates}; 1959 FED; 1963 Yaundé; 1971 SPG; 1975 Lomé; 1992 Maastricht; 1993 Copenhague; 1995 Madrid y Barcelona; 1997 \autor{Pape}; 2000 Cotonú, \autor{Burnside} y \autor{Dollar}; 2001 Todo menos armas; 2002 \autor{Collier} y \autor{Dollar}, \autor{Manners}; 2005 \autor{Alesina}, \autor{Angeloni} y \autor{Etro}; 2004 vecindad, \autor{Nye}, \autor{Schimmelfennig} y \autor{Sedelmeier}; 2007 \autor{Hufbauer} y otros; 2008 UpM; 2009 Asociación Oriental, Lisboa; 2012 \autor{Bradford}; 2017 Consenso; 2019 \autor{Campos} y otros; 2020 metodología, Equipo Europa; 2021 IVDCI, Global Gateway; 2022 tope al petróleo; 2023 Samoa, seguridad económica; 2025 Pacto por el Mediterráneo, Europa Global, Sevilla, «momento del euro global»; 2026 Ucrania y Moldavia abren capítulos.
\end{itemize}

\bloque{Datos clave}
\begin{itemize}
\item UE 2025: 17,9\,\% del PIB mundial; 13,8\,\% de las exportaciones de bienes; 22,7\,\% de las de servicios. Euro: 20,2\,\% de las reservas.
\item Rúbrica 6: 113.700 M€; IVDCI 79.462; IPA III 14.162; humanitaria 11.569.
\item AOD 2025: CAD 0,26\,\%; miembros UE 0,42\,\%; España 0,27\,\%.
\item Global Gateway: 306.000 M€ (2021-2024). Europa Global 2028-2034: $\approx$200.000 M€.
\item Rusia: 21 paquetes; tope 44,1 \$; $\approx$210.000 M€ inmovilizados.
\end{itemize}

\bloque{Preguntas del tribunal}
\begin{enumerate}
\item \textbf{¿Por qué es paralela la competencia de desarrollo?} La UE tiene su política sin impedir la de los Estados (art. 4.4); hay que coordinar 28.
\item \textbf{¿Qué es un acuerdo mixto?} El que toca competencias nacionales (inversión de cartera, ISDS): lo ratifican también los 27 (Dictamen 2/15).
\item \textbf{¿Funciona la condicionalidad?} Si la recompensa es creíble y cercana; fatiga, vetos y retrocesos tras entrar.
\item \textbf{¿Tope o embargo?} El tope mantiene el crudo en el mercado y quita renta; su punto débil, la flota en la sombra.
\item \textbf{¿Quién representa a la UE en el FMI?} No es miembro: el BCE es observador; posiciones del euro por el art. 138 TFUE; EURIMF.
\item \textbf{¿Es AOD el SPG?} No: no transfiere recursos; es cooperación.
\item \textbf{¿Qué es Global Gateway?} Inversión movilizada (300.000 M€) con garantías del FEDS+; crítica: no es dinero desembolsado ni adicional.
\item \textbf{¿Qué cambia con Europa Global 2028-2034?} Un instrumento con IPA y humanitaria; menos foco en la pobreza (90\,\% AOD, sin suelos).
\item \textbf{¿Qué es el Acuerdo de Samoa?} El sucesor de Cotonú con la OEACP: base común y 3 protocolos; provisional desde 2024.
\item \textbf{¿Cuánto da España?} 0,27\,\% de la RNB (2025); 0,7\,\% en 2030 por ley \vertema{4.A.26}.
\end{enumerate}

\begin{trampa}
\begin{itemize}
\item AOD: \emph{4} requisitos; donación equivalente desde los datos de \emph{2018}.
\item Copenhague \emph{1993}; Madrid \emph{1995}; metodología \emph{2020}: 6 bloques.
\item IVDCI: \emph{79.462} M€; incluye el antiguo FED desde \emph{2021}.
\item Samoa: firmado en \emph{2023}, aplicado provisionalmente desde \emph{2024}, \emph{no} en vigor.
\item El tope: 60 \$ (2022) $\to$ \emph{47,6} (2025) $\to$ \emph{44,1} (2026).
\item La UE es \emph{miembro} de la OMC, la FAO y el G-20; \emph{observadora} en la ONU; en el FMI, el BCE es observador.
\end{itemize}
\end{trampa}

\end{columnas}
\end{document}
